% ═══════════════════════════════════════════════════════════════
%  appendices/a-source.tex
%  Приложение А. Исходный код
% ═══════════════════════════════════════════════════════════════

\chapter{Исходный код}
\label{app:sources}

Всё, что описано в этой книге,~--- \emph{работает}. Не <<в теории>>,
не <<в примерах>>, а в реальном приложении \textbf{DualBoard STM32
Oscilloscope}, которое можно собрать и запустить.

\section{Где взять код}

Весь исходный код проекта доступен на GitHub:

\begin{center}
\url{https://github.com/danayer/dualboard-stm32-oscilloscope}
\end{center}

\begin{notebox}
\textbf{Ссылка-заглушка.} Замените URL на реальный адрес вашего
репозитория. Если репозиторий закрытый~--- используйте GitLab,
Bitbucket или собственный сервер.
\end{notebox}

\section{Структура репозитория}

\begin{verbatim}
dualboard-stm32-oscilloscope/
├── firmware/
│   ├── oscilloscope/       # Прошивка осциллографа
│   └── generator/          # Прошивка генератора
├── pc-app/                 # GTK4 приложение
│   ├── main.c
│   ├── Makefile
│   └── ...
├── docs/
│   ├── protocol.md         # Протокол обмена
│   └── ...
└── README.md
\end{verbatim}

\section{Лицензия}

Проект распространяется под лицензией \textbf{MIT}. Это значит:

\begin{itemize}
    \item Можно \emph{использовать} код в любых целях.
    \item Можно \emph{модифицировать}.
    \item Можно \emph{распространять}.
    \item \emph{Обязательно} сохранить упоминание авторства.
\end{itemize}

\begin{ideabox}
\textbf{Почему MIT?} Это самая простая и либеральная лицензия.
Она не накладывает ограничений~--- вы можете делать с кодом
что угодно. Единственное требование~--- сохранить копию лицензии
и упоминание автора.
\end{ideabox}

\section{Как читать исходники}

Код в книге~--- это \emph{объяснение}. Исходники~--- это
\emph{истина}. Если что-то не сходится~--- верьте исходникам.

\begin{enumerate}
    \item Начните с \texttt{main.c}~--- он покажет порядок
          инициализации.
    \item Затем~--- \texttt{app\_data.h}~--- это структура
          всего состояния.
    \item Дальше~--- по главам книги: \texttt{serial\_port.c},
          \texttt{line\_buffer.c}, \texttt{plot\_data.c},
          \texttt{config.c}, \texttt{callbacks.c}, \texttt{ui.c}.
\end{enumerate}

\section{Что дальше}

Если вы нашли баг~--- создайте Issue на GitHub.
Если хотите улучшить~--- сделайте Pull Request.
Если просто используете~--- расскажите, что получилось.
Это лучшая благодарность за труд.
